\chapter{Requisitos do sistema}
\label{ap:requisitos}

Este apêndice reproduz, de forma consolidada, os requisitos funcionais (FR) e
não funcionais (NFR) do projeto, agrupados por etapa de desenvolvimento. Os
requisitos foram escritos antes da implementação e funcionaram como contrato
de cada iteração; cada um possui critério de aceitação associado, no formato
DADO/QUANDO/ENTÃO (ver \autoref{ap:testes} para a suíte correspondente).

\section*{Requisitos funcionais --- \poc{1} (fundação jogável)}
\label{ap:fr1}
\begin{itemize}
  \item \textbf{FR-001} --- O projeto usa um mapa de entrada com ações
  semânticas (mover, olhar, interagir, cancelar, correr, lanterna, inventário,
  pausa, avançar diálogo). É proibido fixar índices de botão ou eixo de
  joystick no \emph{gameplay}.
  \item \textbf{FR-002} --- O jogador se movimenta por teclado (WASD/setas) e
  pelo analógico esquerdo do joystick.
  \item \textbf{FR-003} --- A câmera é controlada por mouse e, quando o
  adaptador expõe os eixos, pelo analógico direito.
  \item \textbf{FR-004} --- Gravidade e colisão impedem atravessar paredes e
  piso.
  \item \textbf{FR-005} --- Menu de pausa acionado por teclado e joystick.
  \item \textbf{FR-006} --- HUD mínimo com objetivo e convite de interação.
  \item \textbf{FR-007} --- Tela de diagnóstico de entrada lista dispositivos,
  eixos e botões em tempo real, com zona morta configurável (padrão 0,15).
  \item \textbf{FR-008} --- Conexão e desconexão a quente do joystick sem
  quebrar a execução.
\end{itemize}

\section*{Requisitos funcionais --- \poc{2} (laboratório funcional)}
\label{ap:fr2}
\begin{itemize}
  \item \textbf{FR-009} --- Sistema de interação desacoplado, com detecção por
  proximidade, convite contextual e sinal de interação.
  \item \textbf{FR-010} --- Inventário diegético (kit técnico) com itens
  coletáveis (fusível F-17, módulo RS-404, chave).
  \item \textbf{FR-011} --- Portas abrem e fecham, podendo depender de
  permissões ou estado.
  \item \textbf{FR-012} --- Equipamentos possuem estados serializáveis
  (desligado, falha, pronto, operando).
  \item \textbf{FR-013} --- Missão principal em seis passos: localizar
  fusível, localizar módulo, instalar ambos, reiniciar CLP, diagnosticar bomba
  e abrir a porta do corredor.
  \item \textbf{FR-014} --- Retroalimentação visual e sonora para interações e
  mudanças de estado.
  \item \textbf{FR-015} --- Sala 3D modelada no Blender e exportada em GLB,
  com identidade visual conforme o projeto de arte do laboratório.
\end{itemize}

\section*{Requisitos funcionais --- \poc{3} (IA local)}
\label{ap:fr3}
\begin{itemize}
  \item \textbf{FR-016} --- Adaptador local com endpoint de conversa que
  valida a entrada, injeta contexto, chama o modelo e valida a saída.
  \item \textbf{FR-017} --- A saída devolve texto e intenção estruturada; o
  jogo aceita apenas as intenções cadastradas.
  \item \textbf{FR-018} --- Cliente assíncrono no jogo com limite de tempo,
  tratamento de erro e \emph{fallback} para mensagens pré-definidas.
  \item \textbf{FR-019} --- Terminal de IA na interface, com entrada de texto,
  histórico e resposta.
  \item \textbf{FR-020} --- A personalidade da IA fica em arquivo de
  configuração versionado, fora dos scripts.
  \item \textbf{FR-021} --- A IA não controla física nem executa comandos;
  devolve apenas texto e intenções que o jogo valida.
\end{itemize}

\section*{Requisitos funcionais --- \poc{4} (mundo reativo)}
\label{ap:fr4}
\begin{itemize}
  \item \textbf{FR-022} --- NPCs reagem ao estado do laboratório.
  \item \textbf{FR-023} --- Memória separada em quatro camadas: sessão, estado
  persistente, conhecimento narrativo e histórico conversacional.
  \item \textbf{FR-024} --- Eventos alteram o mundo: luzes conforme energia,
  portas conforme permissão e tom da IA após alerta ignorado; a iluminação é
  em camadas (luz geral alinhada às luminárias visíveis e acentos por zona); e
  a câmera de segurança acompanha o jogador após a energia ser restabelecida.
\end{itemize}

\section*{Requisitos funcionais --- \poc{5} (laboratório procedural)}
\label{ap:fr5}
\begin{itemize}
  \item \textbf{FR-025} --- Módulos tipados (corredor, controle, bombas,
  baterias, sensores, arquivo, contenção) com entradas, saídas, pontos de
  interesse e iluminação.
  \item \textbf{FR-026} --- Geração determinística por semente.
  \item \textbf{FR-027} --- Conectividade garantida: nenhuma sala inacessível
  e todos os pontos de interesse alcançáveis.
\end{itemize}

\section*{Requisitos funcionais --- \poc{6} (recorte vertical)}
\label{ap:fr6}
\begin{itemize}
  \item \textbf{FR-028} --- Experiência jogável de 10 a 20 minutos, com começo,
  meio e fim.
  \item \textbf{FR-029} --- Tutorial de controles dentro do jogo.
  \item \textbf{FR-030} --- Salvamento e carregamento de progresso.
  \item \textbf{FR-031} --- Áudio ambiente, alarmes e efeitos.
  \item \textbf{FR-032} --- Epílogo variável conforme a escolha do jogador.
  \item \textbf{FR-033} --- Empacotamento em executável e demonstração sem
  intervenção do desenvolvedor.
\end{itemize}

\section*{Requisitos não funcionais}
\label{ap:nfr}
\begin{itemize}
  \item \textbf{NFR-001} --- Executar em computador sem GPU dedicada
  (medido: 60 quadros por segundo limitados por sincronismo vertical em GPU
  integrada, 1280$\times$720).
  \item \textbf{NFR-002} --- Escala do Blender ao motor: 1 unidade = 1 metro.
  \item \textbf{NFR-003} --- Modelo de linguagem local pequeno (alvo
  \texttt{llama3.2:3b}; validação com \texttt{llama3.1:8b}).
  \item \textbf{NFR-004} --- Limite de tempo de 30 segundos na chamada ao
  modelo, com \emph{fallback}.
  \item \textbf{NFR-005} --- Robustez de entrada: nenhum índice de botão fixo
  no \emph{gameplay}.
  \item \textbf{NFR-006} --- Segurança da IA: apenas intenções cadastradas;
  resposta inválida converte para \texttt{NONE}.
  \item \textbf{NFR-007} --- Assets neutros, ajustáveis por pós-processamento.
  \item \textbf{NFR-008} --- Determinismo do laboratório procedural: mesma
  semente, mesma disposição.
\end{itemize}
